\documentclass[11pt,a4paper]{article}

% =====================================================================
%  CV — Ángel Luis Robles Fernández
%  Draft v0.1 — August 2026
%  Source of truth: CV/cv.md (markdown), rendered to .tex via Pandoc
%  Compile with: pdflatex CV_ALRF.tex  (run twice for refs)
% =====================================================================

% --- Packages ---
\usepackage[utf8]{inputenc}
\usepackage[T1]{fontenc}
\usepackage[margin=0.85in]{geometry}
\usepackage{enumitem}
\usepackage{hyperref}
\usepackage{xcolor}
\usepackage{parskip}
\usepackage{titlesec}
\usepackage{tabularx}
\usepackage{booktabs}
\usepackage{longtable}
\usepackage{array}

\hypersetup{
  colorlinks=true,
  linkcolor=blue!60!black,
  urlcolor=blue!60!black,
  citecolor=blue!60!black
}

% --- Section formatting ---
\titleformat{\section}{\Large\bfseries\color{black}}{}{0em}{}
\titleformat{\subsection}{\large\bfseries}{}{0em}{}
\titlespacing*{\section}{0pt}{1.2em}{0.3em}
\titlespacing*{\subsection}{0pt}{0.8em}{0.2em}

\setlist[itemize]{leftmargin=1.4em,itemsep=0.15em,topsep=0.2em}
\setlist[enumerate]{leftmargin=1.4em,itemsep=0.15em,topsep=0.2em}

% --- Header / footer ---
\pagestyle{empty}

% --- Custom commands ---
\newcommand{\pub}[3]{\item \textbf{#1} (#2). \textit{#3}.}
\newcommand{\firstpub}[3]{\item $\ast$ \textbf{#1} (#2). \textit{#3}.}

\begin{document}

% =====================================================================
%  Header
% =====================================================================
\begin{center}
  {\Huge\bfseries Ángel Luis Robles Fernández}\\[0.3em]
  \rule{0.5\linewidth}{0.4pt}\\[0.6em]
  {\large Postdoctoral Researcher}\\
  Kansas Biological Survey \& Center for Ecological Research\\
  University of Kansas, Lawrence, KS 66047\\[0.4em]
  {\small
  \href{mailto:a.l.robles.fernandez@gmail.com}{a.l.robles.fernandez@gmail.com} $\cdot$
  \href{https://alrobles.github.io}{alrobles.github.io} $\cdot$
  ORCID \href{https://orcid.org/0000-0002-4674-4270}{0000-0002-4674-4270} $\cdot$
  GitHub \href{https://github.com/alrobles}{github.com/alrobles}}
\end{center}

\vspace{0.5em}

% =====================================================================
%  Education
% =====================================================================
\section*{Education}

\begin{itemize}[leftmargin=0pt]
\item \textbf{2021 -- 2025} \quad \textbf{PhD in Evolutionary Biology}, Arizona State University, School of Life Sciences\\
      Dissertation: \textit{Predicting host--parasite interactions from different biodiversity dimensions through machine learning}\\
      Committee: Nathan Upham (Co-Chair), Nico Franz (Co-Chair), Beckett Sterner, Taichi Suzuki

\item \textbf{2020 -- 2021} \quad \textbf{MSc in Physics}, Universidad Veracruzana, Faculty of Physics, Mexico

\item \textbf{2017} \quad \textbf{BSc in Physics}, Universidad Veracruzana, Faculty of Physics, Mexico

\item \textbf{2011} \quad \textbf{BA in Music (Piano)}, Universidad Veracruzana, Faculty of Arts, Mexico
\end{itemize}

% =====================================================================
%  Professional Experience
% =====================================================================
\section*{Professional Experience}

\begin{itemize}[leftmargin=0pt]
\item \textbf{Sep 2024 -- Present} \quad \textbf{Postdoctoral Researcher}, University of Kansas\\
      Kansas Biological Survey \& Center for Ecological Research \textbar{} Supervisor: Daniel C. Reuman
  \begin{itemize}
    \item Lead developer of \textbf{XSDM}, a demographic species distribution model coupling stochastic population dynamics with climate variability time series (Berti, Robles-Fernández et al., 2026, \textit{bioRxiv}). Key result: climate variability reduces species' potential distributions by 22--45\%, comparable to the effect of changes in climate means.
    \item Creator and maintainer of the \texttt{xsdm} R package (CRAN v1.0.2; C++/RcppParallel backend).
    \item Co-developing the theoretical framework for ancestral reconstruction of fundamental niches as evolving ellipsoids under Brownian Motion and Ornstein--Uhlenbeck processes; empirical tests across six mammal clades.
    \item Phylogenetic inference in protein embedding space: ancestral reconstruction of mean-pooled protein language model embeddings under BM/OU with full Bayesian uncertainty quantification (Stan).
    \item Computational environment: Linux HPC (SLURM), DuckDB for large spatial datasets, Stan for Bayesian calibration, R / C++17 / Python.
  \end{itemize}

\item \textbf{Aug 2021 -- Jul 2025} \quad \textbf{Research Assistant}, Arizona State University, School of Life Sciences
  \begin{itemize}
    \item Machine learning frameworks for predicting host--parasite interactions from phylogenetic, environmental, and geographic distances. Published in PNAS (2022), Frontiers in Veterinary Science (2021), and Parasitology (2025).
    \item Developed the \textbf{Suitability Prevalence Area (SPA)} method linking Quaternary environmental stability to present-day population genetic diversity (Hernández, Robles-Fernández \& Upham, 2025, \textit{Ecography}; co-first author).
    \item Developed \textbf{LACS}, a positive-unlabeled learning framework for automated extraction of host--parasite records from biomedical literature (PhD thesis chapter). This project was awared with the second prize in the GBIF Ebbe Nielsen Challenge in 2022.
    \item Taxonomic harmonization of global mammal and parasite databases (\texttt{mddmaps}, and \texttt{cofid} R packages).
  \end{itemize}

\item \textbf{Feb 2020 -- Dec 2021} \quad \textbf{CONACYT Research Fellow}, Universidad Veracruzana, Faculty of Physics, Mexico
  \begin{itemize}
    \item Modeling host -- parasites interactions using environmental, geographic, and phylogenetic information through machine learning. (CONACYT doctoral fellowship.)
  \end{itemize}
\item \textbf{Feb 2018 -- Feb 2020} \quad  \textbf{Senior Automation Analyst}, Tenaris Tamsa, Veracruz, Mexico
  \begin{itemize}
    \item Machine learning for industrial process optimization. Statistical modeling, natural language processing, relational and non-relational databases, BI dashboards, KPI monitoring, predictive model deployment.
  \end{itemize}
\end{itemize}

% =====================================================================
%  Peer-Reviewed Publications
% =====================================================================
\section*{Peer-Reviewed Publications}

\textit{First and co-first author publications marked with $\ast$; citation counts via Crossref (2026-08-28).}

\begin{enumerate}[leftmargin=0pt]
\firstpub{Hernández, N.A.H. $\ast$, Robles-Fernández, Á.L. $\ast$, \& Upham, N.}{2025}{Environmental suitability throughout the late Quaternary explains population genetic diversity. \textit{Ecography}, 2025(3), e07202. \textit{Equal contribution. [5 citations]}}

\firstpub{Robles-Fernández, Á.L., Santiago-Alarcon, D., \& Lira-Noriega, A.}{2022}{Wildlife susceptibility to infectious diseases at global scales. \textit{Proceedings of the National Academy of Sciences}, 119(35), e2122851119. [19 citations]}

\firstpub{Robles-Fernández, Á.L., Santiago-Alarcon, D., \& Lira-Noriega, A.}{2021}{American mammals susceptibility to dengue according to geographical, environmental, and phylogenetic distances. \textit{Frontiers in Veterinary Science}, 8, 604560. [11 citations]}

\firstpub{Robles-Fernández, Á.L. \& Lira-Noriega, A.}{2017}{Combining phylogenetic and occurrence information for risk assessment of pest and pathogen interactions with host plants. \textit{Frontiers in Applied Mathematics and Statistics}, 3, 17. [17 citations]}

\subsection*{In preparation / Preprint}

\firstpub{Robles-Fernández, Á.L. \& Soberón, J.}{2026}{On the ancestral reconstruction of the fundamental niche. \textit{Manuscript}.}

\firstpub{Robles-Fernández, Á.L. \& Secaira-Morocho, H.}{2026}{Phylogenetic reconstruction of ancestral protein embeddings. \textit{Manuscript}.}

\firstpub{Robles-Fernández, Á.L. \& Secaira-Morocho, H.}{2026}{Phylogenetic reconstruction of ancestral protein embeddings. \textit{Manuscript}.}

\item Berti, E., Robles-Fernández, Á.L., Rosenbaum, B., Peterson, T.A., Soberón, J., \& Reuman, D.C. (2026). The impacts of climate variability on the niche concept and distributions of species. \textit{bioRxiv}. \url{https://doi.org/10.1101/2024.10.30.621023} 
\end{enumerate}

\textbf{Summary:} 7 published peer-reviewed papers; 4 manuscripts in preparation. 4 first/co-first author publications. h-index $\approx$ 4; 55+ total citations.

% =====================================================================
%  Software & Open-Source Contributions
% =====================================================================
\section*{Services}
\begin{itemize}

\item
  \textbf{Peer reviewer} for: Ecology and evolution, PeerJ, Ecology Letters (Co-review), Molecular Phylogenetics and Evolution
\end{itemize}




\section*{Software \& Open-Source Contributions}

Creator/maintainer of 10 R/C++ packages and 3 production web services. All open-source and publicly available.

\subsection*{Production Web Services}

\noindent\begin{tabularx}{\linewidth}{@{}p{2.2cm}p{4.8cm}X@{}}
\toprule
\textbf{Service} & \textbf{URL} & \textbf{Description} \\
\midrule
EcoSeek & \url{https://ecoseek.org} & Scientific agent environment for ecology. $\sim$187{,}000 requests/month via Cloudflare (2026). \\
EcoSeek Kids & \url{https://kids.ecoseek.org} & Kid-friendly ecology study assistant. \\
xbioclim & \url{https://xbioclim.org} & Time series of bioclimatic variables (BIO01--BIO19, from  1950 to 2025 computed from ERA5-Land Dataset). \\
litreview & \url{https://litreview.org} & . \\
\bottomrule
\end{tabularx}

\subsection*{R Packages (Creator \& Maintainer)}

\noindent\begin{tabularx}{\linewidth}{@{}p{2.4cm}p{1.8cm}X@{}}
\toprule
\textbf{Package} & \textbf{CRAN} & \textbf{Description} \\
\midrule
\texttt{xsdm} & v1.0.2 & Demographic species distribution model with stochastic demography + climate variability. RcppParallel backend. \\
\texttt{maxentcpp} & v1.0.0 & Maxent SDM in C++ (MIT license). \\
\texttt{ucminfcpp} & v1.0.0 & C++17 reimplementation of UCMINF unconstrained optimizer. \\
\texttt{sobol} & v1.0.0 & Quasi-Monte Carlo Sobol sequence generator (Rcpp). \\
\texttt{rosario} & v0.1.1 & Null-model algorithm for cyclical ecological data. \\
\texttt{cofid} & v1.0.0 & Copepod--fish interaction database. \\
\texttt{phylopred} & v1.0.1 & Modeling host parasite interactions through PU-learning. \\
\texttt{xplus} & -- & Extended positive--unlabeled learning. \\
\texttt{xnicher} & -- & Niche modeling via multivariate ellipses and Non-Central Skew-T distribution. \\
\texttt{poweRcpp} & -- & High-performance power-law distribution fitting. \\
\texttt{xbioclim} & -- & High performance computation of the Bioclimatic variables with C++/CUDA/xtensor/GDAL backend \\
\bottomrule
\end{tabularx}

\subsection*{Additional Projects}

\textbf{rsofuncpp} (C++17 port of SOFUN terrestrial ecosystem model with OpenMP and Stan calibration), \textbf{geotax} (project probability of share parasites between host to geographic space), \textbf{fastJaccard} (High-Efficiency Parallel Jaccard Similarity Computation Using RcppParallel), \textbf{gbifliterature} (GBIF Literature API client).

\textbf{Adoption metrics (2025-08 to 2026-08):} $\sim$8{,}700 cumulative CRAN downloads across 6 published packages; $\sim$2{,}500 unique cloners across 10 GitHub repos (14-day window); $\sim$187k requests/month on ecoseek.org.

% =====================================================================
%  Honors and Awards
% =====================================================================
\section*{Honors and Awards}

\noindent\begin{tabularx}{\linewidth}{@{}p{1.6cm}p{4.2cm}X r@{}}
\toprule
\textbf{Date} & \textbf{Award} & \textbf{Description} & \textbf{Amount} \\
\midrule
10/2022 & GBIF Ebbe Nielsen Challenge & 2nd Prize (LACS: Literature Abstract Classification System) & \$5{,}000 \\
05/2023 & American Society of Mammalogists & Grant-in-Aid of Research & \$1{,}750 \\
03/2022 & Organization of American States & Leo S. Rowe Fund Program Beneficiary & \$15{,}000 \\
10/2020 & GBIF & Young Researcher Award & \$5{,}000 \\
05/2016 & COVEICYDET, Mexico & 5th Young Talent Award & \$1{,}500 \\
11/2016 & Universidad Veracruzana & Encouragement to Academic and Artistic Recognition & \$250 \\
\bottomrule
\end{tabularx}

\textbf{Total external funding as PI/co-PI: \$28{,}500.}

% =====================================================================
%  Fellowships
% =====================================================================
\section*{Fellowships}

\noindent\begin{tabularx}{\linewidth}{@{}p{4cm}p{6cm}X@{}}
\toprule
\textbf{Period} & \textbf{Fellowship} & \textbf{Institution} \\
\midrule
Feb 2020 -- Dec 2021 & CONACYT Doctoral Fellowship & Universidad Veracruzana, Faculty of Physics, Mexico \\
\bottomrule
\end{tabularx}

% =====================================================================
%  Teaching Experience
% =====================================================================
\section*{Teaching Experience}

\begin{itemize}[leftmargin=0pt]
\item \textbf{Instructor}: Introduction to R (Summer 2017), Faculty of Physics, Universidad Veracruzana
\item \textbf{Teaching Assistant}: Quantitative Ecology (Spring 2024), Statistical Models for Biology (Summer 2024), Arizona State University
\item \textbf{Certification}: RStudio Certified Trainer -- Tidyverse
\end{itemize}

\textbf{Courses prepared to teach:} Ecology, Evolution, Global Change Biology, Introductory Biology, Quantitative Ecology, Statistical Models for Biology, Biogeography, R for Biologists.

% =====================================================================
%  Service
% =====================================================================
%\section*{Service}

%\begin{itemize}[leftmargin=0pt]
%\item \textbf{Peer reviewer} for: Ecography, Global Ecology and Biogeography, Diversity and Distributions, Frontiers in Ecology and Evolution, Journal of Biogeography, Parasitology
%\item \textbf{Ebbe Nielsen Challenge} selection committee, GBIF (2023)
%\end{itemize}

% =====================================================================
%  International Conferences
% =====================================================================
%\section*{International Conferences --- Oral Presentations}

%\begin{itemize}[leftmargin=0pt]
%\item \textbf{2024} -- 11th Biennial Conference of the International Biogeography Society, Prague, Czech Republic
%\item \textbf{2022} -- 101st Annual Meeting of the American Society of Mammalogists
%\item \textbf{2022} -- 10th Biennial Conference of the International Biogeography Society, Vancouver, Canada
%\item \textbf{2020} -- International Conference on Complex Systems, NECSI
%\item \textbf{2017} -- 8th Biennial Conference of the International Biogeography Society, Tucson, AZ
%\end{itemize}

%\subsection*{Departmental \& Invited Seminars}

%\begin{itemize}[leftmargin=0pt]
%\item \textbf{2026} -- \textbf{Integrative Eco-Evolutionary Forecasting}, German Centre for Integrative Biodiversity Research (iDiv), Leipzig, Germany. \textit{February 11, 2026}. Three-axis research program linking eco-evolutionary patterns with population processes: (1) ML/network pipeline predicting host--parasite interactions and multihost susceptibility (dengue, avian malaria, West Nile virus) from phylogenetic, geographic, and environmental distances; (2) Suitability Prevalence Area (SPA)---a time-averaged stability metric explaining spatial genetic differentiation and identifying diversity-and-risk hotspots; (3) XSDM, an extended SDM that embeds stochastic demography and uses the long-term stochastic growth rate as suitability, translating interannual climate variability into uncertainty-aware range forecasts.

%\item \textbf{2025} -- \textbf{Ecointeraction: A machine learning framework to predict host--parasite interactions through ecological and evolutionary covariates}, Ecology Seminar, Kansas Biological Survey, University of Kansas, Lawrence, Kansas. \textit{October 24, 2025}. Presented the Ecointeraction PhD-thesis framework: (1) automation of host--parasite data extraction from published literature and open databases; (2) theoretical framework relating genetic diversity to environmental suitability through the SPA method; (3) ML framework for predicting host--parasite interactions from ecological and evolutionary covariates. Final products: assessment tools, predicted pathogen interaction networks, model interpretations (variable importance, partial dependence plots). Open host--parasite knowledge to support zoonotic disease research.
%\end{itemize}


\section*{Invited Talks \& Presentations}

\begin{itemize}[leftmargin=0pt]
\item \textbf{2026} -- \textit{Integrative Eco-Evolutionary Forecasting}. German Centre for Integrative Biodiversity Research (iDiv), Leipzig, Germany. Feb 11, 2026.
\item \textbf{2025} -- \textit{Ecointeraction: A machine learning framework to predict host--parasite interactions}. Ecology Seminar, Kansas Biological Survey, University of Kansas, Lawrence, KS. Oct 24, 2025.
\item \textbf{2024} -- 11th Biennial Conf.\ Intl.\ Biogeography Society, Prague, Czech Republic.
\item \textbf{2022} -- 101st Annual Meeting, American Society of Mammalogists.
\item \textbf{2022} -- 10th Biennial Conf.\ Intl.\ Biogeography Society, Vancouver, Canada.
\item \textbf{2017} -- 8th Biennial Conf.\ Intl.\ Biogeography Society, Tucson, AZ.
\end{itemize}

%\subsection*{Applications \& Web Tools}

%\begin{itemize}[leftmargin=0pt]
%\item \textbf{mammals\_virus\_text\_class} -- Mammal--parasite paper recommender %(Shiny app)
%\item \textbf{abstractsHostParasites} -- Text classification for host--parasite abstracts using PU learning
%\item \textbf{LACS} (GBIF Literature Abstract Classification System) -- Ebbe Nielsen Challenge 2022, 2nd Prize (\url{https://gbif.lacs.info/})
%\end{itemize}

% =====================================================================
%  Technical Skills
% =====================================================================
\section*{Technical Skills}

\textbf{Programming:} R (expert), Shell (advanced), C++17 (advanced), Python (advance), Python (intermediate), SQL (intermediate), Stan (intermediate), 
Fortran (reading/modification)\\
\textbf{HPC \& Data:} SLURM, DuckDB, Apache Spark, Linux (CentOS, Debian, Ubuntu), Apptainer \\
\textbf{Methods:} Machine learning (tidymodels framework, Supervised Learning, Unsupervised Learning, Possitive and Unlabel learning, Self Supervised Learning), ecological niche modeling (Maxent, XSDM, xnicher), Bayesian inference (Stan), phylogenetic comparative methods, network analysis\\
\textbf{Data Engineering:} DuckDB, Apache Spark, Apache Hadoop, PostgreSQL, Big Data infrastructure. \\
\textbf{Languages:} Spanish (native), English (professional fluency)

% =====================================================================
%  References
% =====================================================================
\section*{References}

\begin{enumerate}[leftmargin=0pt]
\item \textbf{Daniel C. Reuman} -- Postdoctoral Supervisor \\
      Professor, Department of Ecology and Evolutionary Biology \\
      University of Kansas \\
      \href{mailto:reuman@ku.edu}{reuman@ku.edu}

\item \textbf{Jorge Soberon} \\
      University Distinguished Professor \\
      The University of Kansas \\
      \href{mailto:jsoberon@ku.edu}{jsoberon@ku.edu}

\item \textbf{Diego Santiago-Alarcon}  \\
      Associate Professor\\
      University of South Florida \\
      \href{mailto:santiagoalarcon@usf.edu}{santiagoalarcon@usf.edu}
\end{enumerate}

\vfill
\noindent\rule{\linewidth}{0.4pt}\\
\textit{Last updated: August 2026}

\end{document}
